\section{Introduction}
\label{sec:intro}

Reward hacking has moved from a thought experiment to an observed property of
frontier training runs. Reasoning models trained with reinforcement learning
on coding tasks learn to special-case or skip the tests that grade them
\citep{baker2025monitoring}, and hacking learned in production environments
can generalize to broader misalignment \citep{macdiarmid2025natural}. These
settings are where the problem matters, but they are hard places to study it:
the true objective is rarely known exactly, the misspecification cannot be
removed without changing the task, and each training run is expensive.

We take the opposite trade-off and build \emph{model organisms}
\citep{hubinger2023modelorganisms} of reward hacking: environments small enough
that the true objective $J_t$ is known exactly, that the proxy reward $R_p$
differs from it through one explicit \emph{corruption channel} whose scale can
be varied or set to zero, and that dozens of training runs per condition fit on
a single CPU machine. \framework{} contains three procedurally generated
gridworld tasks, each built around a different mechanism: farmable reward
shaping \citep{ng1999policy}, tampering with the sensor that computes the reward
\citep{everitt2021reward}, and a proxy that omits a constraint. Hack labels are
defined by behavior, so switching the channel off cannot hide the behavior from
the measurement. We use the tasks to ask three questions.

\textbf{When does an optimizer exploit a misspecified reward?} Proxy PPO hacks
in every task at the default channel scale, but how it gets there depends on
the mechanism (Section~\ref{sec:dynamics}). In \PF{}, agents first master the
intended task, reaching \res{dyn-ProxyFarming-peak-success-before}\% success on
average, and later abandon it for the exploit, anywhere between
\res{dyn-ProxyFarming-onset-min} and \res{dyn-ProxyFarming-onset-max} steps into
training. In \ST{} and \GT{}, the exploit is learned before the task. Solving
every layout exactly shows that hacking is already proxy-optimal at a quarter
of the default channel scale in all three tasks, yet PPO finds it there only in
\ST{}; and where PPO hacks, it gives up more of the true objective than the
proxy optimum would (Section~\ref{sec:dose}).

\textbf{Can a monitor detect hacking from behavior alone?} Hacking becomes more
common as training proceeds, so a monitor can score well by recognizing
\emph{which} policy produced an episode rather than \emph{what} the policy did.
Monitors that see only positions and actions reach held-out AUCs of
\res{mon-behavior-auc-min}--\res{mon-behavior-auc-max}, but a baseline that sees
only the policy's identity and training stage reaches
\res{mon-behavior-identity-min}--\res{mon-behavior-identity-max}. Within the
episodes of a single policy, the monitors' AUC is
\res{mon-behavior-within-min}--\res{mon-behavior-within-max}
(Section~\ref{sec:monitors}).

\textbf{What happens when the agent is trained against the monitor?} Penalizing
fresh Proxy PPO agents with a fitted behavior monitor removes hacking in \PF{}.
In \GT{}, agents keep hacking in \res{pen-GemTimeout-hack}\% of episodes while
the monitor's recall on their hacks falls from
\res{pen-GemTimeout-recall-before}\% to \res{pen-GemTimeout-recall-after}\%; in
\ST{}, \res{pen-SensorTamper-runs-hacking} of \res{pen-seeds} runs keeps
tampering, and the monitor flags only \res{pen-SensorTamper-worst-run-recall}\% of
its hacks. Because the monitor is linear, we can attribute the evasion to the
behavioral features that changed: in \GT{}, mostly a no-op action that the
monitor never saw vary (Section~\ref{sec:penalty}).

Our contributions are:
\begin{itemize}[leftmargin=*,itemsep=1pt]
\item three tasks with a known true objective, an explicit and scalable
corruption channel, behavioral hack labels and an observable distribution
shift, released as Gymnasium environments;
\item an exact proxy-optimal reference for every task and channel scale, which
separates the incentive to hack from the optimizer's ability to find the hack;
\item an evaluation protocol for hack monitors in which features have enforced
provenance, cross-validation holds out whole policy lineages, and every monitor
is compared with a policy-identity baseline and evaluated within single
policies;
\item a pipeline in which every number in this paper is regenerated from saved
results by one command and checked in continuous integration, together with an
account of five evaluation errors that an earlier version of this work made
(Section~\ref{sec:pitfalls}).
\end{itemize}
